\scanpage{432}
\mspage{344}

Soit $S$ une partie de $\mathrm{Ob}\,\unsure{M|\mathcal M}$, $\underline S$ la sous-catégorie pleine qu'elle définit. On dit que $S$ est une \uline{partie} $\unsure{\Phi|\mathscr I|\mathscr D}$-$\pi$-\uline{accessible} de $\mathrm{Ob}\,M$, ou que $\underline S$ est une \uline{sous-catégorie} $\Phi$-$\pi$-\uline{accessible} \uline{dans} $M$, si \unsure{,} les conditions suivantes sont satisfaites

\margin{\unsure{condition|conditions} prél. sur $M$ $\leftarrow$}

$1^\circ)$ $\pi$ est utile pour $M$, $M$ est stable par $\varinjlim$ de type $\Phi$.

$2^\circ)$ $\underline S$ est une sous-catégorie $\pi$-accessible \unsure{:} dans $M$.

$3^\circ)$ $\underline S$ est \deleted{accessible} stable dans $M$ par $\varinjlim$ de type $\Phi$.

% slanted
\margin{\uline{NB} Quand $\Phi=\{I_{\alpha_0}\}$, on dit aussi que $\underline S$ est $\alpha_0$-$\pi$-accessible dans $M$. Si $\Phi$ est formé des cat. de cardinal $\leq\pi$ et grdes devant $\pi_0$, on dit que $\underline S$ est $\pi_0$-$\pi$-\unsure{accessible} dans $M$.}

\uline{NB} Quand la condition préliminaire $1^\circ)$ \uline{sur $M$} est satisfaite, \inserted{la validité} des conditions $2^\circ)$ $3^\circ)$ \unsure{équivaut|équivalent} donc à la condition que $\underline S$ soit \uline{strictement} pleine dans $M$, et que \inserted{de plus} $\underline S$ satisfasse la condition a) du th. 4.13.22. D'autre part, \unsure{cette|ces} condition est équivalente à \deleted{aux} \inserted{chacune des} \unsure{conditions} a', b, c de ce même théorème, et aussi à la dernière condition c'), \unsure{sous} réserve qu'on ait $\pi^{\pi_0}=\pi$ (pour un cardinal $\pi_0<\pi$ utile pour $M$) et que $J\in\Phi$ implique que $J$ est \unsure{grde|grd} devant $\pi_0$.

Conceptuellement, c'est la définition donnée dans 4.13.23, \deleted{qui \unsure{p}} i.e. la condition \textcircled{a} de 4.13.22, qui paraît \unsure{le|la} plus \unsure{clair|claire}, ou sa variante immédiate a'\deleted{\unsure{a)|)}}. Mais du point de vue pratique, du \deleted{\unsure{le mani}} \unsure{travail} avec les parties accessibles, c'est la propriété \textcircled{c} qui paraît le plus commode. Le poids principal du th. 4.13.22 est bien

% archive sheet number: 431

